\cotefile{157-3}
\scanpage{282}
\mspage{232}

une $W$-cofibration, $X'_0\unsure{-|\to}X'$ est dans $W$, et pour la $=$ raison $Y'_0-Y'$ l'est aussi. Donc $f'$ l'est aussi.

\uline{NB}\textsuperscript{\uline{bis}} C'est un argument formel, \inserted{(disons)} valable \inserted{(\unsure{sous} forme duale)} dans tout $(M,W)$ où on a un \unsure{théorème} de factorisation pour une flèche quelconque $f:X-Y$ en $X\xrightarrow{i_f}\overline X\xrightarrow{\alpha}Y$, avec $i_f\in W$ et $\alpha$ une $W$-fibration. Alors, si $X,Y$ sont $W$-fibrés sur $S$, \deleted{\unsure{toute|tout}} \unsure{une|e} $S$-flèche de $X$ dans $Y$ qui est dans $W$, est dans $W_S^{\unsure{\sim}}$. Et dualement, si \deleted{\unsure{toute}} flèche \inserted{$X-Y$} se factorise en $X'\longrightarrow\deleted{\unsure{X}}\widetilde Y\longrightarrow Y$, avec $X\longrightarrow\widetilde Y$ une $W$-cofibration et $\widetilde Y-Y$ \deleted{dans} $W$.

Il me vient ici une \inserted{(petite)} inquiétude : ai-je au moins prouvé que si
\[
S\longrightarrow X
\]
est une $W$-cofibration, alors pour toute flèche $S'\longrightarrow S$ dans $Cat$, le carré cocartésien dans $Cat$
\diagram{\begin{tikzcd}[ampersand replacement=\&]
 S \arrow[r,dash] \arrow[d] \& X \arrow[d] \\
 S' \arrow[r] \& X'
\end{tikzcd}}
est aussi $W$-\deleted{\unsure{Cat-cartésien|cocartésianique|cocartésien}} \inserted{cocartésien} \inserted{dans le sens ancien, spécial à $Cat$ (déf. 3 p. \unsure{131|191}) Pour le comparer à la théorie générale, cf prop 20 p. 262} ? \deleted{\unsure{Autrement|Effectivement},}

% archive sheet number: 282

